\section{Introduction}
\label{sec:intro}

Garment manipulation is a contact-rich, bimanual task that remains
difficult to automate directly. Learning-based approaches, such as
vision-language-action policies, depend on high-quality demonstration
data. That data in turn depends on a teleoperation interface that is
precise, low-latency, and comfortable enough for long collection
sessions. This paper documents such an interface: a teleoperation system
that drives two SO-101 arms (the LeRobot/SO-ARM100
lineage~\citep{soarm100}) from Meta Quest hand controllers, with a MuJoCo
digital twin used both for operator rehearsal and for quantitative
benchmarking of design alternatives.

Three properties of the SO-101 shape every design decision. First, each
arm has only five degrees of freedom (shoulder pan, shoulder lift, elbow
flex, wrist flex, wrist roll). There is no wrist-yaw joint, so full 6-D
pose tracking is structurally over-constrained: a naive 6-D tracker
trades position error against an orientation error it can never
eliminate. Second, the arm is small (\(\approx\)\SI{0.41}{m} maximal
reach from the shoulder-lift pivot) while a human operator's natural
gestures easily exceed that reach, so out-of-envelope hand targets are
the norm, not the exception. Third, the STS3215 bus servos are
position-controlled with modest stiffness, and this makes command
smoothness a first-class concern. Each servo runs its own internal
position loop towards the latest goal angle, so the commanded trajectory
is effectively a sequence of step references. Jerky commands produce
large instantaneous steps, and a low-stiffness gear train answers large
steps with overshoot and audible oscillation rather than crisp motion.
Repeated large reversals also load the plastic gears and the serial bus
itself, so rough commands cost hardware life as well as tracking
fidelity. Smoothness therefore cannot be left to the servos; the
teleoperation pipeline must deliver it.

This document makes five contributions:
\begin{enumerate}
  \item a description of the \emph{armplane} retargeting pipeline:
        grip-clutched absolute-position mapping, input smoothing with
        the One-Euro filter (an adaptive low-pass for interactive
        input, defined in Section~\ref{sec:prelim}), and the armplane
        orientation construction, which builds only orientations the
        5-DoF arm can actually realise (Section~\ref{sec:method});
  \item an analytic \emph{workspace envelope} for the SO-101 with four
        selectable out-of-envelope target policies~--- warn, project,
        freeze, and slowdown~--- applied uniformly to every
        teleoperation input source (Section~\ref{sec:envelope});
  \item the integration of a \emph{Telegrip-style split IK}
        \citep{telegrip} as a pluggable method: position-only damped
        least squares on the three proximal joints with a direct
        analytic wrist mapping, targeting wrist agility
        (Section~\ref{sec:telegrip}; the full port, and how it departs
        from the upstream implementation, is specified in
        Appendix~\ref{app:telegrip});
  \item a simulation benchmark comparing six IK methods on tracking,
        wrist-agility, and out-of-envelope suites, with metrics that
        include orientation error and roll phase lag
        (Section~\ref{sec:experiments});
  \item a within-subjects \emph{user study protocol} (\(n\approx 5\))
        pairing the objective benchmark with subjective evaluation on a
        bimanual coordination task (Section~\ref{sec:exp-userstudy}),
        using NASA-TLX (the NASA Task Load Index, a six-scale
        self-reported workload questionnaire~\citep{nasatlx}), SUS (the
        System Usability Scale, a ten-item usability
        questionnaire~\citep{sus}), custom teleoperation-feel scales,
        and semi-structured interviews.
\end{enumerate}

The rest of the paper reads as follows. Section~\ref{sec:related}
positions the system against other bimanual and low-DoF teleoperation
work. Section~\ref{sec:prelim} introduces the common anatomy of
headset-to-arm teleoperation for readers new to the area.
Section~\ref{sec:system} describes the hardware and software.
Section~\ref{sec:method} presents the retargeting method, including what
happens to the gripper orientation at the moment teleoperation engages.
Section~\ref{sec:experiments} reports the benchmark and the user-study
protocol, and Section~\ref{sec:conclusion} concludes.

The paper is a \emph{living document}: it is the canonical reference for
the teleoperation side of the repository and is updated in the same
branch as any change to the methods, mappings, envelope, or benchmark
results.
